% Generated draft: requires mathematical and source review.
\documentclass[11pt]{article}
\usepackage[margin=25mm]{geometry}
\usepackage{fontspec}
\setmainfont{Latin Modern Roman}
\usepackage{amsmath,amssymb,mathtools,unicode-math}
\setmathfont{Latin Modern Math}
\usepackage{enumitem,needspace,xurl,hyperref}
\setlength{\parindent}{0pt}
\setlength{\parskip}{5pt}
\setlength{\emergencystretch}{3em}
\newcommand{\N}{\mathbb{N}}
\newcommand{\Z}{\mathbb{Z}}
\newcommand{\Q}{\mathbb{Q}}
\newcommand{\R}{\mathbb{R}}
\newcommand{\C}{\mathbb{C}}
\newcommand{\F}{\mathbb{F}}
\newcommand{\A}{\mathbb{A}}
\newcommand{\PP}{\mathbb{P}}
\newcommand{\E}{\mathbb{E}}
\newcommand{\eps}{\varepsilon}
\newcommand{\dd}{\,\mathrm{d}}
\begin{document}
\section*{UnsolvedMath ID: 6200103}
\textbf{Source problem code:} catalogue entry
\subsection{Definitions and mathematical statement}
An $\mathbb R$-tree is a geodesic metric space with a unique geodesic arc between any two points. The source asks to extend Rips' theory of group actions on trees to higher-dimensional building settings, including products of $\mathbb R$-trees. It does not specify the acting groups, stability/smallness hypotheses, or which structural conclusions of Rips' theory should generalize.

For a CAT(0) metric space $X$, use the rank definition supplied in the item: $\operatorname{rank}(X)\geq n$ if every geodesic segment is contained in a subset $E\subseteq X$ isometric to a nondegenerate flat Euclidean $n$-dimensional parallelepiped. For a locally CAT(0) space $Y$ with the indicated universal cover, set $\operatorname{rank}(Y)\geq n$ when its universal cover has this property.

Develop an appropriate higher-dimensional structure theory and a rank-rigidity analogue beyond Riemannian manifolds. The model result quoted in the source is that a compact nonpositively curved Riemannian manifold of rank at least two has a nontrivial Riemannian product splitting of its universal cover or is locally symmetric. The extracted sentence also contains an unattached finite-cover clause; it does not clearly state a separate global product conclusion for a finite cover of the manifold. Its exact wording should be checked rather than strengthened.

For general CAT(0) or locally CAT(0) spaces, the source supplies neither a complete list of extra hypotheses nor a precise decomposition/building/symmetry conclusion. Such a theorem must specify its category, completeness conditions, and rank-rigidity alternatives; no universal classification of all rank-at-least-two CAT(0) spaces is asserted by this draft.
\subsection{Short English statement}
Extend Rips-type structure theory to higher-dimensional building actions, and formulate rank rigidity for general CAT(0) spaces using the source's flat-parallelepiped rank condition. The general hypotheses and conclusions are not fully specified.
\subsection{Sources}
\begin{itemize}
\item Misha Kapovich, \emph{Problems on Boundaries of Groups and Kleinian Groups} (2005), Problem 103, PDF page 25, \url{https://www.math.ucdavis.edu/~kapovich/EPR/problems.pdf}; accessed 2026-07-29, PDF-text-normalized extraction.
\item Ballmann, source reference [1], and Burns--Spatzier, source references [15,16], are the rank-rigidity attributions in the extract. The record supplies no corresponding bibliography entries.
\item The supplied summary mentions Dymarz, Hume, and Rips complexes on buildings without publication identifiers or an exact statement connecting those constructions to the requested Rips theory.
\item Local registry \texttt{data/unsolvedmath/problems.json}, permanent UnsolvedMath ID $6200103$, source code catalogue entry, \emph{Boundaries of Groups and Kleinian Groups --- Problem 103}. Dataset attribution: Ulam AI and the UnsolvedMath Contributors.
\end{itemize}
\end{document}
